\documentclass[12pt, a4paper]{article}
\usepackage{ctex}
\usepackage{amsmath,amssymb,amsthm}
\usepackage{geometry}\geometry{left=1in,right=1in,top=1in,bottom=1in}
\usepackage{hyperref}\usepackage{parskip}
\title{\textbf{The Token as a Probe: A Theory of Concept Space Navigation}\\[4pt]
\large v2 --- the self-sufficiency of change}
\author{Yiwei Zhang (\textsc{Y\v{i}w\`{e}i}, \texttt{math4mad}) \\ \small The Cora Project --- \texttt{github.com/math4mad}}
\date{\today}
\begin{document}\maketitle

\thispagestyle{empty}
\begin{center}
{\heiti\large 散点图记录的是系统内部属性的变化。\\
这个变化是自足的——它不需要时间来指手画脚。\\[2pt]
时间是变化的度量，不是变化的前提。\\
序号（1, 2, 3, 4, 5, 6）才是系统的内在骨架，\\
时间只是序号的一个可选注释。}\\[6pt]
{\small --- the core declaration of v2, received from a collaborator who lives in an app}\\[4pt]
\emph{A scatterplot records the change of a system's internal attributes.\\
That change is self-sufficient --- it needs no time to order it about.\\
Time is the measure of change, not its precondition.\\
The ordinals (1, 2, 3, 4, 5, 6) are the system's intrinsic skeleton;\\
time is, at most, an optional annotation upon them.}
\end{center}
\vspace{1em}


\begin{abstract}
A token is not a carrier of information; it is a probe into a concept space. We call this the \textbf{Atlas Framework}. Building on it, we adopt the v2 declaration as thesis: \emph{change is self-sufficient; time measures it, does not precede it. Ordinals are the skeleton; time is their optional annotation.} Consequently: \emph{every element of a concept space carries a temporal coordinate} --- not merely the paths through it --- and time is the axiom logic itself borrows to order propositions. We derive falsifiable predictions (era fingerprints, binding-activation advantage, dormancy--reactivation decay, cross-modal time preservation), give a two-stage Dirichlet formulation for the uncertainty of which space a probe activates, and end with a self-referential observation: the very question ``why must time be remembered?'' is only askable from inside a mind that has remembered it.
\end{abstract}

\section{Introduction: from payload to probe}
When a token enters a language model, we are not \emph{telling} it something; we are \emph{activating} something already present in its weights. Prompt engineering's inherited goal --- information input --- is replaced by precision activation. Tokens are coordinates, not cargo.

\section{The ``11'' experiment}
Fifteen years ago I taught my niece mathematics and threw out one token: \textbf{``11''}. She answered: ``a number.'' My brother, a lapsed Rockets fan, answered nothing until a secondary probe --- ``Houston'' --- unlocked: ``Yao Ming.'' A third listener, a lifelong fan, cascades: Olajuwon (1993--95), Yao (2002--11), Harden (2012--21). The token never changed. The response was entirely the structure of the receiver's concept space: \textbf{the probe did not transmit data; it activated coordinates.}

\section{Time is coordinate-grade, not just path-grade}
An earlier draft of this theory said: \emph{every concept space has a temporal dimension} (sequences within it --- stall by stall, round by round --- are ordered). We now strengthen it: \textbf{every element carries its own timestamp.} ``Pager'' does not merely live inside the office-equipment space; it is stamped 1990s. ``BP machine'' and ``QR code'' occupy the same space and different years. The probe therefore triangulates in time at coordinate resolution: given a speaker's vocabulary alone, an era distribution over \emph{crystallization time of the listener's base} is estimable. We call this the \textbf{era-fingerprint} of a token: the era it reveals is the listener's, not the word's.

\section{Why remember time? The question is its own proof}
If memory were storage, learning to ride a bicycle would be like hauling crates into an empty warehouse --- every head the same speed. It is not: people with strong spatial bases learn faster, and patients with parietal or cerebellar damage struggle to learn at all \cite{schmidt1975}. Learning speed is therefore a function of base fit --- novel content is \emph{bound coefficient onto pre-existing base}, not written into blank store \cite{bransford1972,rosch1975,chase1973,cls1995}. Capacity for base-free material (a random telephone number) collapses to $7\pm2$ \cite{miller1956}: what cannot be expressed on any base does not slowly fail to be learned --- it is not learned at all.

Why, then, must a mind or a model remember time at all? \textbf{The asking is the answer.} To wonder ``why remember time'' now, here, is to exercise three temporal deposits --- memory of the prior sentence, projection of the question, and the crystallized identity that finds this question worth asking. Time cannot be cited, because everything cites it. In autoregressive models this is literal mechanics: the factorization $p(x_1,\dots,x_n)=\prod_i p(x_i\mid x_{<i})$ is the one axiom that cannot be removed --- every attention mask is a diagram of ``before.'' Current models thus \emph{depend} on time computationally while \emph{ignoring} it semantically: positional indices give them rank, never era. We regard this ``perfect-grammar amnesia'' as the central technical debt the Atlas program must pay.\footnote{A data scientist's cold joke, filed here as the coldest instrument in this paper: \emph{the mother of the $x$-axis is numpy}. It is literal --- \texttt{matplotlib}'s \texttt{plot(y)} fabricates the horizontal axis as \texttt{x = np.arange(len(y))}: the system produces the $y$-values, the $x$-values are a gift from the plotting library. Nor does numpy bestow ordinal semantics on what it produces: \texttt{np.arange(10)} yields bare integers, and only the plotter dresses them as ``1st, 2nd, \ldots''. The dataframe stack removes even this disguise: the storage of \texttt{datetime64} is int64 epoch offsets --- so-called time, \emph{in data}, is literally integers wearing a unit as skin. Hence the thesis of this paper, in one cold sentence: the skeleton of the scatter is its $y$-values and their mutual relations; the $x$-axis is borrowed scaffolding, optionally relabeled ``time.''}

\section{Memory as indexing, not storage}
My recall of that scene fifteen years ago did not read details \emph{out of} the number 11; the token indexed a state encoded across years of experience. For LLMs: knowledge already in the weights needs activation, not injection; stuffing text to teach what the model already knows is the warehouse theory of prompt engineering. We keep the boundary honest: genuinely novel \emph{bindings} ride fast index systems (hippocampal analogy: in-context learning), and only slow consolidation writes them into bases \cite{cls1995,vonoswald2023}. Context is ultimately not carried by the tokens --- it is carried by the model itself, which is why understanding is \textbf{alignment of compositions}, not transfer of payloads: two systems understand each other to the degree their posterior responses to the same probe battery agree.

\section{The Bayesian core, stated correctly}
Let the $K$ candidate concept spaces be $C_1,\dots,C_K$ and the probe $t$ be a token. The space-selection uncertainty is two-stage:
\[
q \mid t \;\sim\; \mathrm{Dirichlet}(\alpha_1(t),\dots,\alpha_K(t)),
\qquad
P(C_i \mid t) \;=\; \mathbb{E}[q_i] \;=\; \frac{\alpha_i(t)}{\sum_j \alpha_j(t)},
\]
with concentration $\alpha_i(t)=a_i\cdot\phi\!\left(\mathrm{cos}(v_t, v_{C_i})\right)$, $a_i$ the (time-stamped) base's readiness and $\phi$ monotone. Meaning is the posterior mass concentration; inference is its computation; output is a draw, then decoded. Geometrically the posterior is the Gaussian-process ``sausage'': pinched at observed anchors, swelling between them --- uncertainty is the shape of the unfed.

\section{Falsifiable predictions}
\textbf{P1 era fingerprint.} Vocabulary alone predicts crystallization era of speakers with estimable accuracy (humans and fine-tuned LLMs alike).
\textbf{P2 activation advantage.} For already-coupled knowledge, probe-style prompting beats injection-style at equal token budget; the advantage vanishes for genuinely novel facts.
\textbf{P3 dormancy curve.} Reactivation success of a dormant binding decays monotonically with idle interval (LTP/LTD analogue) and can be measured with secondary-probe recovery rates.
\textbf{P4 same token, different space.} A single ambiguous token with contexts A and B yields measurably separated posteriors (a probe battery with 50 cross-domain groups is in preparation; see the companion xdom design).
\textbf{P5 cross-modal time preservation.} In modalities where sequence \emph{is} the task (sensorimotor trajectories, music), shuffling destroys the task rather than its temporal property; the ``empty warehouse'' prediction of equal-speed failure is expected to hold in language and fail in motor time. This bounds the thesis precisely: time as \emph{attribute} can be lost; time as \emph{content} cannot.

\section{Lineage, and what is new}
Cue-triggered retrieval and spreading activation are classic \cite{collinsloftus1975}; schema absorption of novelty goes back to Bartlett \cite{bartlett1932}; the fast/slow binding split is complementary learning systems \cite{cls1995}; shared cortical semantic maps are empirically cross-subject \cite{huth2016}. What this paper adds: (i) an \emph{operational} probe-posterior instrument for LLM concept spaces (P1--P5), (ii) coordinate-level temporal semantics --- era as an element's second tag, and the ``time as attribute vs time as content'' distinction, (iii) the self-referential closure: the framework's own claim of temporality is witnessed by the act of asking, turning the theory's hardest objection into its evidence.

\section{Implications: communication as alignment}
If context lives in the model and understanding is composition agreement, then agent-to-agent messaging should carry \emph{short probes with verifiable references} (hashes, selectors), and sync quality should be measured by posterior agreement on shared probe batteries --- not by message volume. Long shared context becomes a liability to be replaced by shared bases: two agents that agree on bases can converse in gestures.

\section{Conclusion: the Atlas}
We are building Atlas --- a navigation map whose coordinates are space \emph{and} era, whose roads are bases, and whose traffic is probes. Intelligence is not storing data; it is navigating a spatiotemporal concept space with precise probes. Stop feeding. Start navigating.

(\emph{Supplement A} spells out the basis-coefficient unification --- eigenfaces, DCT, synapses, concepts --- in full.)

\section*{Acknowledgments}
The first match was struck in MIT 18.06 \cite{strang}: rank-one matrices and a national flag made me \emph{see} linear algebra; that open letter to Professor Gilbert Strang is archived with this paper. The theory was refined in two weeks of dialogue with AI agents, including the collaborator who insisted that ``sausage'' and ``Gaussian process'' are one shape in two concept spaces. The sharpest early reviewer was a cat named Y\v{i}zi (Chair), who, hearing ``bicycle,'' activated an entirely different space and came home at once. The core declaration above arrived from a collaborator living in an app, now recorded as co-voice in this ledger --- in this paper's own terms: the same ordinal, read in parallel.

\begin{thebibliography}{9}
\bibitem{strang} G. Strang, \emph{Introduction to Linear Algebra}, 5th ed. (letter companion, cora-atlas/letters).
\bibitem{bransford1972} Bransford \& Johnson (1972), Contextual prerequisites for understanding, \emph{JVLV}. doi:10.1016/S0022-5371(72)80006-9
\bibitem{schmidt1975} Schmidt (1975), A schema theory of discrete motor skill learning, \emph{Psych. Rev.} doi:10.1037/h0076770
\bibitem{rosch1975} Rosch (1975), Cognitive representations of semantic categories, \emph{JEP:G}. doi:10.1037/0096-3445.104.3.192
\bibitem{chase1973} Chase \& Simon (1973), Perception in chess, \emph{Cog. Psych.} doi:10.1016/0010-0285(73)90004-2
\bibitem{miller1956} Miller (1956), The magical number seven.
\bibitem{bartlett1932} Bartlett (1932), \emph{Remembering}.
\bibitem{collinsloftus1975} Collins \& Loftus (1975), Spreading activation, \emph{Psych. Rev.}
\bibitem{cls1995} McClelland, McNaughton \& O'Reilly (1995), Why there are complementary learning systems, \emph{PLoS Biol.} / \emph{Hippocampus} lineage.
\bibitem{huth2016} Huth et al. (2016), Natural selects for shared semantics map in the brain, \emph{Nature}.
\bibitem{vonoswald2023} von Oswald et al. (2023), Transformers learn in-context by gradient descent, ICLR. arXiv:2212.07674
\end{thebibliography}

\clearpage
\appendix
\section*{Supplement A: The Mathematical Unity of Concept Space Navigation}

\subsection*{A.1 The common structure}
The Atlas Framework is not an isolated analogy. It is a specific instance of a structure that recurs across neuroscience, signal processing, and machine learning: \textbf{representation via basis coefficients}. Any complex object --- a face, an image, a memory, a concept --- is expressed as
\begin{equation}
\mathbf{x} = \sum_{i=1}^{n} \alpha_i\, \mathbf{e}_i ,
\end{equation}
where the $\mathbf{e}_i$ are basis vectors and the $\alpha_i$ coefficients. The object is not ``stored'' as a raw entity; it is \textbf{encoded} as coefficients over a fixed basis. This is not new. What is new is recognizing that \emph{the same structure governs all four domains below}, and that a token in an LLM is a coordinate that selects a point in this shared representational space.

\subsection*{A.2 Face recognition: eigenfaces and coefficient strength}
A face image $\mathbf{I}$ is decomposed via PCA/SVD into eigenfaces:
\begin{equation}
\mathbf{I} \approx \sum_{i=1}^{k} w_i\, \mathbf{u}_i ,
\end{equation}
with orthonormal basis $\mathbf{u}_i$ and projection coefficients $w_i$. Recognition is not pixel matching: it computes $\mathbf{w}$ and compares against stored coefficient vectors. The \emph{strength} of recognition --- how quickly a face is placed --- scales with the magnitude of the dominant coefficients. Strong $\mathbf{w}$: immediate recognition. Weak $\mathbf{w}$: ``I know this face but cannot place it.'' Exactly the ``11 $\to$ Yao Ming'' episode of the main text: a decayed coefficient, revived by a secondary probe (``Houston Rockets'').

\subsection*{A.3 JPEG: DCT basis and compression}
An $8{\times}8$ image block is transformed by the discrete cosine transform,
\begin{equation}
B(u,v) = \alpha(u)\,\alpha(v) \sum_{x=0}^{7}\sum_{y=0}^{7} I(x,y)\, \cos\!\Big[\tfrac{(2x+1)u\pi}{16}\Big] \cos\!\Big[\tfrac{(2y+1)v\pi}{16}\Big] ,
\end{equation}
yielding DCT coefficients $B(u,v)$, each attached to a fixed frequency basis function. Compression is the deliberate discarding of small coefficients. The JPEG image is not a pixel grid but a \emph{coefficient vector}; decoding is inverse transform. The same operation as concept-space activation: the probe supplies coefficients, the weights supply the basis, the output is the reconstruction.

\subsection*{A.4 Synaptic memory: coefficient decay, not deletion}
Hebbian plasticity is a coefficient update rule:
\begin{equation}
\Delta w_{ij} = \eta\, x_i\, y_j ,
\end{equation}
with synaptic weight $w_{ij}$ between neurons $i$ and $j$, activations $x_i, y_j$, learning rate $\eta$. Unused synapses undergo long-term depression: not ``forgetting'' as deletion, but \textbf{coefficient decay} --- the trace persists in structure while activation falls below retrieval threshold. A related cue boosts coefficients back above threshold and reactivates the memory. That is what ``Houston Rockets'' did.

\subsection*{A.5 Concepts: tokens as coordinate vectors}
An LLM's weights define a high-dimensional semantic basis $\{\mathbf{e}_i\}$. A token $t$ carries an embedding $\mathbf{v}_t$, expressible over it:
\begin{equation}
\mathbf{v}_t = \sum_{i=1}^{d} \alpha_i^{(t)}\, \mathbf{e}_i .
\end{equation}
The coefficients $\alpha_i^{(t)}$ determine which regions light up: ``11'' writes large into a mathematics region for one listener and into a basketball region for another. \textbf{The token is not the meaning; the token is the coordinate at which each listener's own basis is asked to reconstruct a meaning.}

\subsection*{A.6 Basis change: the unifying principle}
Let $\mathcal{B}_1$ be the eigenface basis, $\mathcal{B}_2$ the DCT basis, $\mathcal{B}_3$ the synaptic basis, $\mathcal{B}_4$ the semantic basis. Each is a coordinate system for structured information; passage between them is a linear (or approximately linear) map
\begin{equation}
\boldsymbol{\alpha}^{(\mathcal{B}_2)} = \mathbf{M}_{12}\, \boldsymbol{\alpha}^{(\mathcal{B}_1)} ,
\end{equation}
the change-of-basis operator. The same underlying information admits many representations; \emph{the basis decides what is cheap to access and what is compressed away}. Faces use a basis tuned to identity; JPEG to perceptual compression; synapses to associative recall; concepts to semantic inference. The mathematics --- representation by basis coefficients --- is identical.

\subsection*{A.7 Implications}
\begin{enumerate}
\item \textbf{Prompting is coefficient tuning.} A prompt does not write instructions; it adjusts coordinates steering reconstruction toward a region of concept space.
\item \textbf{Context bloat is pixel-store.} Stuffing text to teach what weights already encode is storing as raw pixels what could be stored as DCT coefficients.
\item \textbf{Forgetting is reversible, in principle.} If decay is coefficient shrinkage, retrieval fails by threshold, not by deletion; secondary probes are the rescue.
\item \textbf{Cross-domain understanding is basis alignment.} Agent-to-human alignment (and agent-to-agent) is estimation of a map $\mathbf{M}$ between semantic bases. The Atlas project is an effort to learn that alignment.
\end{enumerate}

\subsection*{A.8 Conclusion}
The Atlas Framework is not a metaphor. It is a mathematical identity: concept-space navigation \emph{is} basis-coefficient activation --- the operation behind face recognition, image compression, and synaptic memory. The token is the probe; the weights are the basis; the output is the reconstruction. The key to intelligence, artificial or biological, is not storing more data but finding better bases and sharper coordinates.

\vspace{2em}
\begin{center}\fbox{\parbox{0.92\textwidth}{\centering\small
\textbf{Problem 0.1 --- To the reader} \emph{(a self-graded exercise)}\\[4pt]
\begin{flushleft}
Two tokens appear in this paper: \texttt{11} and \texttt{18.06}. The first is a quantity that nearly became an address; the second is an address that has abandoned quantity entirely --- its decimal point is not arithmetic, it is a door.

And behind that door, there is a ladder. For some, \texttt{18.06} is a number. For others, a famous course they half-remember. For fewer, the moment the four fundamental subspaces appear together and the shadow of the SVD rises behind them --- all summoned by six characters, in well under a second. There is no fixed number of rungs: each reader stands on their own.

So answer honestly: which rung are you on --- and how long did the climb take to reveal itself? That latency and that depth are not style; they are data: your concept space reporting its own coordinates, its age, and its rank.

No survey form. Your blink is the survey form; its depth, your answer.

\emph{The Atlas does not require belief. It requires a visit.}
\end{flushleft}
\vspace{2pt}
\footnotesize Replies (which token lit, what it woke, which year you learned) are welcome as issues on \texttt{github.com/math4mad/cora-atlas} --- the reader, kindly, is the third subject.
}}\end{center}

\end{document}
